\vspace*{-3pt}
\section{Conclusion and Future Work}
\vspace*{-3pt}
In this paper, for the first time, we show that PHB with large learning rate induces large catapults, resulting in a much larger sharpness reduction than that of GD.
We first provide empirical evidence for this on linear diagonal networks (LDNs) and nonlinear neural networks.
We then hypothesize that the large catapult of PHB is caused by momentum {\it prolonging} self-stabilization~\citep{damian2023stabilization}.
We verify our hypothesis for ReLU scalar networks and LDNs and rigorously prove that it holds.

% as well as a scalar two-layer ReLU network, and offer an intuition of why such \emph{large catapults} occur from a self-stabilization perspective~\cite{damian2023stabilization}: momentum prolongs the sharpness reduction.

% Lastly, we empirically validate this phenomenon in 
\noindent\textbf{Future Work.}~~
Firstly for PHB, we sometimes observe a lack of PS after the large catapults, as shown in Figure~\ref{fig:cifar10-exp}(a--c). However, with larger datasets, PS appears to occur again, as shown in Figure~\ref{fig:cifar10-exp}(d--f). We conjecture that the lack of PS in ReLU networks is due to catapults aggressively reducing the number of active ReLU neurons~\citep{lu2020dying,andriushchenko2023sam}; refer to Appendix~\ref{app:overshooting} for a detailed discussion. Understanding the presence and absence of PS in different setup is an interesting direction.
% A more precise explanation of this phenomenon is left for future works.
% Firstly, for PHB, we sometimes observe a lack of PS after the large catapults, as shown in Figure~\ref{fig:cifar10-exp} (a-c). However, with larger datasets, PS appears to occur again, as shown in Figure~\ref{fig:cifar10-exp}(d-f). In general, the presence or absence of PS after a large catapult seems to depend on the dataset's size and the model architecture. One possible explanation for the lack of PS in {\it ReLU networks} is that the large catapults in PHB aggressively reduce the number of active ReLU neurons~\citep{lu2020dying,andriushchenko2023sam}; see Figure~\ref{fig:deadrelu} of Appendix~\ref{app:dead}. A more precise explanation of this phenomenon is left for future works.
Secondly, although PHB drives iterates towards flatter minima via the large catapults, we do not always observe better generalization.
% this is in contrast to the claims of \cite{zhu2023catapult,zhu2023catapultquadratic}, where the authors claim that catapults lead to better AGOP~\citep{radhakrishnan2022agop}, and thus better generalization.
Further exploring the connection between catapults and generalization is an interesting area for future research and may contribute to the growing literature on the connection between flatness and generalization~\citep{hochreiter1997flat,dinh2017sharp,even2023diagonal,andriushchenko2023sharpness}.

% The hypothesis that flat minima generalize better~\citep{hochreiter1997flat} has been actively studied, with ``counterexamples'' being recently proposed~\citep{dinh2017sharp,even2023diagonal,andriushchenko2023sharpness}. However, no solid conclusion has been reached, and exploring the connection between catapults and generalization would be an interesting area to explore. 

% \begin{remark}[Absence of Progressive Sharpening? - how to bridge catapult and EoS?]
% \label{rmk:dead}
%     For PHB, we sometimes observe a lack of PS after the large catapults, as shown in Figure~\ref{fig:cifar10-exp} (a-c). However, with larger datasets, PS appears to occur again, as shown in Figure~\ref{fig:cifar10-exp}(d-f).
%     In general, the presence or absence of PS after a large catapult seems to depend on the dataset's size and the model architecture.
%     One possible explanation for the lack of PS in ReLU networks is that the large catapults in PHB aggressively reduce the number of active ReLU neurons~\citep{lu2020dying,andriushchenko2023sam}; see Figure~\ref{fig:deadrelu} of Appendix~\ref{app:dead}
% \end{remark}


% \begin{remark}[Flatness and Generalization - future work]
%     We are not claiming that flatter minima~\cite{hochreiter1997flat} necessarily lead to better generalization, and in fact, some ``counterexamples'' have been recently proposed~\citep{dinh2017sharp,even2023diagonal,andriushchenko2023sharpness}. Instead, we solely focus on the effect of momentum and learning rate warmup on how sharpness changes throughout training.
% \end{remark}

